\chapter{Jouer un coup à la souris}

Tous les autres chapitres sont écrits du point de vue de \textbf{celui qui écrit
un moteur}. Celui-ci est écrit du point de vue de \textbf{celui qui joue}. C'est
le point de vue qui manquait, et ce chapitre existe à cause d'une phrase :

\begin{quote}
« Le labo compile bien mon code du palier 1 et un menu déroulant sur la gestion
des ressources de fusion est apparu dans les règles, mais \textbf{il n'y a pas de
bouton pour jouer la fusion}. Et je pense que c'est la même chose pour les
pouvoirs et les artefacts. »
\end{quote}

Cette phrase était juste sur les trois points. Voici, pour chacun, ce qu'on voit
à l'écran et où l'on clique.

\begin{nkretenir}[Ce qui a été vérifié, et comment]
Ce qui est écrit ici sur \emph{dupliquer} a été observé à l'écran : atelier
lancé, un totem cliqué, capture d'écran --- le contour de sélection et les
anneaux verts étiquetés \texttt{D} y sont. La fusion, les pouvoirs et
l'ambiguïté ont été \textbf{mesurés hors interface} par le banc
\nkfile{Applications/ConquerorLab/tests/ambiguite.ps1}, qui appelle directement
la fonction décidant quel coup part au clic : \textbf{13 vérifications, 0
échec}. Ce qui n'a été ni vu ni mesuré est dit comme tel, en clair, plutôt que
déduit du contrat.
\end{nkretenir}

\section{La règle unique : deux clics}

Il n'y a \textbf{qu'un seul geste} dans l'atelier, quel que soit le genre du
coup :

\begin{enumerate}
  \item \textbf{un clic sur une case de départ} --- elle prend un contour
        orange, et toutes les cases où ce totem peut agir s'allument ;
  \item \textbf{un clic sur une case allumée} --- le coup part.
\end{enumerate}

Il n'existe \textbf{aucun bouton} «~Fusionner~», «~Lancer un pouvoir~» ou
«~Ramasser l'artefact~», et il n'en manque pas : ce sont les \textbf{cases} qui
portent l'action.

Un clic ailleurs annule la sélection ; un clic sur un autre de vos totems y
déplace la sélection. \textbf{On ne peut pas jouer un coup illégal} : une case
qui n'est pas allumée ne répond pas. Et si rien ne s'allume, c'est que ce totem
n'a aucun coup --- pas que l'atelier vous ignore.

\section{Ce que les marques veulent dire}

\begin{center}
\begin{tabular}{@{}p{5.2cm}p{8.6cm}@{}}
\toprule
\textbf{Marque} & \textbf{Sens} \\
\midrule
contour orange épais & la case \textbf{sélectionnée} \\
anneau vert $+$ \texttt{D} & destination d'un \textsc{dupliquer} \\
anneau orange plein $+$ \texttt{F3} & la \textbf{case du résultat} d'une \textsc{fusion} à 3 cases \\
anneau orange fin, relié par un trait & une case que la fusion va \textbf{consommer} \\
anneau vert $+$ \texttt{P0}, \texttt{P1} & cible d'un \textsc{pouvoir}, et \textbf{lequel} \\
\texttt{$\times$2} au-dessus d'une case & \textbf{plusieurs coups différents} visent cette case \\
anneaux rouges & les ennemis que le coup \emph{survolé} retournerait \\
\bottomrule
\end{tabular}
\end{center}

Ces marques restent affichées \textbf{tant que la sélection dure}. C'est un point
de conception, pas un détail : ce qui n'apparaît qu'au \emph{survol} n'existe pas
pour qui ne sait pas où passer la souris. Avant le 23 août 2026, la case du
résultat d'une fusion ne se distinguait d'aucune autre destination --- d'où
l'impression qu'il manquait un bouton.


Tous les coups se jouent du même geste : \textbf{un clic sur une case de départ,
un clic sur la case d'arrivée}. Ce qui change d'un genre de coup à l'autre,
c'est \emph{quelle} case est le départ et \emph{quelle} case est l'arrivée.
C'est la seule chose à savoir, et c'est ce que le plateau dessine.

\section{Dupliquer}

Cliquez votre totem : il prend le contour orange, et les cases voisines
atteignables prennent un anneau vert marqué \texttt{D}. Cliquez l'une d'elles :
le coup part.

\section{Fusionner --- où est la « case du résultat »}

Une fusion consomme \emph{plusieurs} cases et n'en produit qu'une. Il n'y a donc
pas de case de départ unique, et c'est ce qui déroute :

\begin{itemize}
  \item cliquez \textbf{n'importe laquelle} des cases que la fusion consomme ---
        elles fonctionnent toutes ;
  \item l'atelier trace alors un trait \textbf{orange} de chaque case consommée
        vers la \textbf{case du résultat}, qui porte un anneau plein et
        l'étiquette \texttt{F} suivie du nombre de cases ;
  \item cliquez cette case du résultat : la fusion part. Le journal affiche
        \textsc{fusionner}, et le totem réapparaît un niveau plus haut.
\end{itemize}

Vous ne sélectionnez donc \textbf{pas} les cases une par une. Le moteur a déjà
énuméré les groupes possibles dans \texttt{CoupsPossibles} ; l'interface ne fait
que retrouver celui que vous désignez.

\begin{nkpiege}[« Il n'y a pas de bouton pour jouer la fusion »]
C'est la phrase exacte d'un stagiaire, et elle décrivait un vrai défaut. Dans
les ateliers antérieurs au \textbf{23 août 2026}, la sélection filtrait sur le
champ \texttt{from} du coup. Or une fusion laisse \texttt{from} \emph{à zéro} et
met ses cases dans \texttt{fuseCells}. Cliquer une case à consommer ne la
sélectionnait donc même pas : aucun contour, aucun anneau, rien --- d'où
l'impression qu'il manquait un bouton. Pire : \texttt{from} mis à zéro
\emph{vaut} la coordonnée $(0,0)$, si bien que le coin du plateau passait pour
la source de toutes les fusions. Si votre atelier se comporte ainsi,
\textbf{votre version est périmée} : reprenez le kit.
\end{nkpiege}

\section{Les pouvoirs, et le menu « Quel coup ? »}

Un pouvoir est désigné par \emph{trois} choses : le lanceur (\texttt{from}), la
cible (\texttt{to}) et \textbf{lequel} (\texttt{powerId}). Deux pouvoirs
différents d'un même totem sur une même cible partagent donc leur départ et leur
arrivée : la case cliquée ne suffit plus à dire quel coup vous voulez.

L'atelier ne choisit pas à votre place. Quand plusieurs coups visent la case que
vous cliquez, il affiche « $\times$2 » \emph{avant} le clic, puis ouvre un menu
\textbf{« Quel coup ? »} qui les nomme --- \texttt{POUVOIR n1},
\texttt{FUSIONNER 2 cases} --- et indique entre parenthèses combien de totems
chacun retourne. Tant que ce menu est ouvert, le plateau ne répond plus : un
clic ne peut pas le traverser et jouer un coup à votre insu.

\begin{nkretenir}
Deux coups qui partagent leur destination restent deux coups \emph{différents},
et vous devez pouvoir jouer les deux. En revanche, deux pouvoirs générés avec le
\textbf{même} \texttt{powerId} sont identiques octet pour octet : l'atelier n'en
montrera qu'un, et ce n'est pas un bug de l'atelier. Numérotez vos pouvoirs.
\end{nkretenir}

Côté écriture, le raccourci \texttt{Pouvoir(joueur, de, vers, idPouvoir)} de
\texttt{ConquerorRegleFacile.h} met le coup à zéro pour vous, comme
\texttt{Dupliquer} et \texttt{Fusionner}.

\section{Les artefacts}

Il n'y a \textbf{pas} de genre de coup « artefact » dans le contrat :
\texttt{NkcMoveKind} vaut \texttt{None}, \texttt{Duplicate}, \texttt{Fuse},
\texttt{Power} ou \texttt{Pass}. Un artefact est un \emph{champ de case},
\texttt{NkcCellView::artefact} --- un identifiant, $-1$ quand il n'y en a pas.
Il se ramasse donc en jouant un coup ordinaire vers la case qui le porte, ou se
déclenche par un \textsc{pouvoir} dont vous choisissez l'identifiant.

C'est un choix de conception, pas un oubli : ajouter un genre de coup casserait
l'ABI de tous les modules déjà compilés, alors qu'un champ de case n'oblige
personne. Aucun moteur livré ne pose d'artefact aujourd'hui --- c'est le palier
2 : le champ existe, il vaut $-1$ partout, et il vous attend.


\section{Où lire la suite}

\begin{itemize}
  \item Écrire un moteur qui \textbf{génère} des fusions : chapitre \emph{Les
        règles en trois fonctions}, section «~Fusionner~», et l'exemple complet
        \nkfile{exemples/rules/RegleFusion.cpp}.
  \item Le contrat nu, champ par champ (\texttt{fuseCells}, \texttt{powerId},
        \texttt{artefact}) : chapitre \emph{Écrire des règles}.
  \item Lire le plateau, les panneaux, le journal et la campagne : chapitre
        \emph{Se servir de l'atelier}.
\end{itemize}

\begin{nkexo}
Reprenez \nkfile{RegleFusion.cpp}, ajoutez \textbf{deux} pouvoirs qui visent la
même case, et vérifiez à l'écran que le menu «~Quel coup ?~» propose bien les
deux et que \textbf{les deux partent}. Si l'un des deux ne part jamais, regardez
d'abord vos \texttt{powerId} : deux coups qui ne diffèrent par rien sont un seul
coup.
\end{nkexo}
